\documentclass[11pt]{article}
\usepackage[margin=1in]{geometry}
\usepackage{amsmath,amssymb,amsthm}
\newtheorem{definition}{Definition}
\newtheorem{theorem}{Theorem}
\newtheorem{remark}{Remark}
\DeclareMathOperator{\spanop}{span}
\DeclareMathOperator{\tr}{tr}
\begin{document}

\section*{T\_E018\_TripleNoGo: Bounded-Grammar Triple-Bridge No-Go}

\subsection*{Status and Grade}

This is a Mode T consolidation theorem for the E018 QM--GR construction thread.
Its grade is:
\[
  \text{finite-carrier diagnostic theorem, bounded to }G^*.
\]
It is not an unconditional analytic theorem and not an absolute impossibility
claim.  It records the construction-grade exit state of the bounded cascade.

\begin{definition}[The bounded grammar \(G^*\)]
Let \(G^*\) be the union of the following audited grammar fragments:
\[
G^*=G_{\mathrm{rec}}\cup G_{\mathrm{atom}}\cup G_{\mathrm{run}}.
\]
\begin{itemize}
\item \(G_{\mathrm{rec}}\): the recognition imports tested in Steps 1--4:
holography/AdS--CFT with RT/HRT, LQG/canonical quantum geometry, and
asymptotic safety.
\item \(G_{\mathrm{atom}}\): the posited-atom-shadow designs \(G_{v1}\)--\(G_{v3}\):
shared currency \(u_i\), route-holonomy \(v\), holographic-RG flow \(w\), and
the bounded triple-audit search over these rows.
\item \(G_{\mathrm{run}}\): the Step 8 local stochastic binary-ring run grammar
\(G_{v4}\), with sector readouts measured after the run and an out-of-sample
interaction witness.
\end{itemize}
\end{definition}

\begin{definition}[Triple shared-completion bridge]
The triple bridge \(B_T\) is the demand for one completion \(E\) and one P6 audit
\(D\) that jointly couple the three sector rows:
\[
  A=\text{area as entanglement shadow price},\quad
  G=\text{amplitude over geometry},\quad
  R=\text{P3 route closure}.
\]
In the finite-carrier diagnostics \(B_T\) is represented by the row \(t\), or
in the run grammar by the demand that the measured sector plateau lies near a
single diagonal fixed point with an independent held-out witness.
\end{definition}

\begin{definition}[Non-circular generation]
A mechanism in \(G^*\) generates \(B_T\) non-circularly only if all of the
following hold:
\begin{enumerate}
\item \(B_T\) is not inserted as a native source row or substrate input.
\item \(B_T\) is not read off from the pairwise bridge union \(u_i,v,w\) alone.
\item In a run grammar, the substrate is not post-processed or defined so that
the three sector readouts are equal by construction.
\item Any witness used for \(B_T\) has independently measured content not
reducible to pairwise statistics and not equal to \(t\).
\end{enumerate}
\end{definition}

\begin{definition}[Rejected smuggling routes]
The three smuggling routes rejected in this thread are:
\begin{enumerate}
\item \emph{Insertion}: use \(t\) itself as a source row or input.
\item \emph{Pairwise union}: treat \(\spanop(u_i,v,w)\) as a triple package.
\item \emph{Rigging}: force the run-readout plateau onto the diagonal, or place
an equivalent triple objective into the substrate.
\end{enumerate}
An orthogonal witness with no audited map to \(t\) is not a smuggling route, but
it is insufficient.
\end{definition}

\begin{theorem}[Bounded triple-bridge no-go in \(G^*\)]
Within \(G^*\), the triple shared-completion bridge \(B_T\) is not closable
non-circularly.  More precisely, no audited mechanism in
\[
G_{\mathrm{rec}}\cup G_{\mathrm{atom}}\cup G_{\mathrm{run}}
\]
generates one shared completion fixed point plus one P6 audit across \(A,G,R\)
without falling into insertion, pairwise-union overread, rigging, or
insufficiency.
\end{theorem}

\begin{proof}
The proof is by audited exhaustion of the bounded grammar \(G^*\).

\paragraph{Recognition fragment \(G_{\mathrm{rec}}\).}
Step 4 recombines the three recognized profiles from holography, LQG, and
asymptotic safety.  The sector-only span has normalized residual \(0\), but this
is the rejected sector-union control.  When compatibility rows are included,
the four bridge rows have normalized residual \(1.0\) and the coupled target has
raw residual \(4.0\), normalized residual \(4/7=0.5714285714285714\), at both
refinement levels \(r=2,4\).  Thus the existing recognition imports do not
supply the triple bridge inside \(G_{\mathrm{rec}}\).

Artifact citation:
\[
\texttt{steps/step4\_mode\_c\_recombination\_artifacts/step4\_schema.json}.
\]

\paragraph{Posited-atom fragment \(G_{\mathrm{atom}}\).}
Steps 5 and 6 generate the three pairwise bridges using \(u_i\), \(v\), and
\(w\), leaving only \(B_T\).  Step 7 then exhausts the bounded \(G_{v3}\)
triple-audit search:
\[
\Xi(B_T\mid \spanop(u_i,v,w))=1.0
\]
at \(r=2\) and \(r=4\), and the coupled residual remains \(1/7\).  Adding a
decimal form, this is \(0.14285714285714285\).  Adding a
dependent pairwise-union row leaves the same residual.  Adding an orthogonal
witness row without an audited map to \(t\) also leaves the same residual.
Adding \(t\) itself gives zero residual, but that is insertion and is excluded
by non-circular generation.

Artifact citations:
\[
\texttt{steps/step7\_mode\_b\_gv3\_triple\_artifacts/step7\_schema.json},
\quad
\texttt{steps/step7\_mode\_b\_gv3\_triple\_artifacts/triple\_audit\_search\_step7.csv}.
\]

\paragraph{Run fragment \(G_{\mathrm{run}}\).}
Step 8 tests a structurally different finite stochastic substrate at \(N=32\)
and \(N=64\).  The measured completion residuals against the diagonal
fixed-point subspace are
\[
0.4947207905172907,\qquad 0.5371910513210945.
\]
The held-out interaction witness gains are
\[
-0.02326224393247789,\qquad -0.08539192789918792.
\]
Both refinements fail the declared run-candidate thresholds.  A forced diagonal
completion has near-zero residual, but it is post-run rigging and is excluded.
The pairwise-union statistic is measured only as a rejected control.

Artifact citations:
\[
\texttt{steps/step8\_mode\_b\_gv4\_run\_triple\_artifacts/step8\_schema.json},
\quad
\texttt{steps/step8\_mode\_b\_gv4\_run\_triple\_artifacts/witness\_audit\_step8.csv}.
\]

The three fragments cover \(G^*\).  Each route either leaves a positive
diagnostic residual or reaches zero only by a rejected smuggling route.  Hence
\(B_T\) is not non-circularly closable within \(G^*\).
\end{proof}

\subsection*{Seven-Gate Audit Summary}

\begin{enumerate}
\item Bounded grammar declared: \(G^*\) is explicitly scoped.
\item Typed target declared: \(B_T\) is the one shared completion and P6 audit.
\item Non-circularity declared: insertion, pairwise union, and rigging are
excluded.
\item Finite-carrier evidence recorded: Steps 4, 7, and 8 provide the residuals.
\item Target-lineage non-equivalence preserved: the theorem is a no-go scoped to
\(G^*\), not a root closure claim.
\item Nonclaim boundary explicit: no claim is made about mechanisms outside
\(G^*\).
\item Next grammar obligation declared: \(G_{v5}\) must supply independent
triple records via a different dynamics class or a named recognition import.
\end{enumerate}

\begin{remark}[Next obligation]
The surviving obligation is
\[
\texttt{G\_E018\_IndependentDynamicsOrRecognition\_v5}.
\]
Literal identifier: \verb|G_E018_IndependentDynamicsOrRecognition_v5|.

It must provide independent triple witness records and pass the same no-smuggling
checks.  Another native copy of \(t\), another pairwise-union statistic, or a
post-run forced diagonal does not meet this obligation.
\end{remark}

\end{document}
